% ════════════════════════════════════════════════════════════════════
%  tầng a, phần 2 — định nghĩa «đủ» (giờ mới đủ ký hiệu), rồi c2 · c5
% ════════════════════════════════════════════════════════════════════

\begin{frame}[fragile]{giờ mới định nghĩa được «đủ» — và nó cần hai điều kiện}
  \framesubtitle{định nghĩa}
  \lead{đã có $v(S)$ ở slide trước và $\mathrm{adq}(S)$ từ \texttt{c0}, nên nói
        được thế nào là \emph{đủ}. điểm mấu chốt: \hl{phải cả hai}.}

  \[
    \boxed{\;
      S \text{ đủ} \;\iff\;
      \underbrace{v(S) \ge \tau_q}_{\text{model trả lời đúng}}
      \;\wedge\;
      \underbrace{\mathrm{adq}(S)}_{\text{các đoạn thật sự nói ra đáp án}}
    \;}
  \]

  \vspace{0.5em}
  \begin{columns}[t]
  \begin{column}{0.46\textwidth}
    \begin{badb}[title={nếu chỉ dùng $v$ — quay lại đoạn 19}]
      \begin{dia}
v({19}) = 0,71          cao
adq({19}) = SAI         phủ 0 / 2 mảnh
      \end{dia}
      \vspace{0.25em}
      đoạn 19 chỉ nhắc \emph{tên phim}; model tự nhớ ra đáp án rồi trả lời đúng.
      \vspace{0.25em}

      bản chỉ-dùng-$v$ sẽ \no{nhận đoạn 19 vào đáp số}, rồi đi dạy bộ truy hồi
      lần sau đi lấy đúng loại đoạn đó. hỏng từ gốc.
    \end{badb}
  \end{column}
  \begin{column}{0.5\textwidth}
    \begin{bx}[title={và lệch này đã được người khác đo}]
      \footnotesize
      \begin{tabular}[t]{@{}>{\raggedright\arraybackslash}p{0.76\linewidth}r@{}}
        \multicolumn{2}{@{}l@{}}{\hl{model trả lời đúng khi ngữ cảnh không đủ}}\\[0.15em]
        \quad trên các bộ đo trong bài & 35--62\,\%\\[0.55em]
        \multicolumn{2}{@{}l@{}}{\hl{trong các ca model trả lời sai}}\\[0.15em]
        \quad do \emph{không có} ngữ cảnh & 10,2\,\%\\
        \quad do ngữ cảnh \emph{có mà không đủ} & 66,1\,\%\\[0.55em]
        \multicolumn{2}{@{}l@{}}{\hl{sau truy hồi, tỉ lệ ca thật sự đủ ngữ cảnh}}\\[0.15em]
        \quad musique-ans & 46,2\,\%\\
        \quad hotpotqa & 44,6\,\%\\
      \end{tabular}
      \vspace{0.35em}

      \dd{joren et al., iclr 2025 — arXiv:2411.06037}
    \end{bx}
  \end{column}
  \end{columns}
\end{frame}

\begin{frame}{hai điều kiện cho bốn ô — và mỗi ô một nguyên nhân khác}
  \framesubtitle{định nghĩa}
  \lead{giữ hai điều kiện \emph{tách rời} thì mỗi ca hỏng rơi vào đúng một ô, và
        ô đó chỉ thẳng ra thủ phạm.}

  \begin{center}
  \begin{tabular}[t]{@{}>{\raggedright\arraybackslash}p{0.12\linewidth}
                        >{\raggedright\arraybackslash}p{0.37\linewidth}
                        >{\raggedright\arraybackslash}p{0.37\linewidth}@{}}
    \toprule
    & \textbf{$\mathrm{adq}$ phủ hết mảnh} & \textbf{$\mathrm{adq}$ không phủ}\\
    \midrule
    \hl{$v \ge \tau_q$}
      & \yes{nhãn tốt} — nhận
      & \no{mồi tham số} — loại, đây là đoạn 19\\[0.7em]
    \hl{$v < \tau_q$}
      & bằng chứng \emph{có} mà model không nối lại được\newline
        $\rightarrow$ lỗi \textbf{model sinh}, tầng b \no{không sửa được}
      & pool \emph{thiếu} đoạn cần thiết\newline
        $\rightarrow$ thoát sớm ngay ở \texttt{c0}\\
    \bottomrule
  \end{tabular}
  \end{center}

  \vspace{0.8em}
  \begin{keyb}
    nếu gộp hai điều kiện thành một điểm số thì mất bảng này, và mất luôn khả năng
    tách \emph{lỗi đi tìm} khỏi \emph{lỗi model đọc}.
    ô dưới-trái là ô mà bản chỉ-dùng-$v$ \no{không bao giờ nhìn thấy}.
  \end{keyb}
\end{frame}

\begin{frame}[fragile]{«nhỏ nhất» nghĩa là gì — ba khái niệm từ lý thuyết trò chơi}
  \framesubtitle{định nghĩa}
  \lead{coi mỗi đoạn là một \emph{người chơi}, mỗi tập đoạn là một \emph{liên minh}.
        $v$ là tiền mà liên minh đó kiếm được. ngôn ngữ này đã có sẵn từ 1965.}

  \begin{columns}[t]
  \begin{column}{0.455\textwidth}
    \begin{bx}[title={ví dụ ba đoạn $a$, $b$, $c$ — dùng lại ở cuối tầng a}]
      \begin{dia}
tau_q = 0,84

v(rỗng)      0,05
v({a})       0,12
v({b})       0,18
v({c})       0,71   <- ĐƠN LẺ CAO NHẤT
v({a,b})     0,94   <- thắng
v({a,c})     0,60   <- thêm a lại GIẢM
v({b,c})     0,66
v({a,b,c})   0,80   <- THUA, vì có c
      \end{dia}
    \end{bx}
  \end{column}
  \begin{column}{0.505\textwidth}
    \begin{bx}[title={ba khái niệm, đọc thẳng từ bảng bên trái}]
      \footnotesize
      \begin{tabular}[t]{@{}l@{\hspace{1em}}>{\raggedright\arraybackslash}p{0.72\linewidth}@{}}
        \hl{thắng} & $v(S) \ge \tau_q$. ở đây chỉ $\{a,b\}$ và không ai khác\\[0.45em]
        \hl{tối tiểu} & thắng, mà \emph{bỏ bất kỳ đoạn nào} cũng thành thua\\[0.45em]
        \hl{mwc} & tập thắng tối tiểu. $W = \{\{a,b\}\}$\\
      \end{tabular}
    \end{bx}
    \begin{badb}[title={$v$ KHÔNG đơn điệu — thêm đoạn có thể làm tệ đi}]
      $v(\{c\}) = 0{,}71$ mà $v(\{a,c\}) = 0{,}60$; và $\{a,b\}$ đang thắng thì
      thêm $c$ làm nó \no{thua}. đoạn gây nhiễu làm model lạc thật, không phải
      nhiễu đo. ta \yes{không sửa} — mwc vốn không cần đơn điệu.
    \end{badb}
  \end{column}
  \end{columns}
\end{frame}

\begin{frame}[fragile]{c2 · một mô hình rẻ đoán điểm thay cho llm}
  \framesubtitle{c2}
  \lead{bước \texttt{c5} sắp tới phải so hàng nghìn tập con. gọi llm cho từng tập
        thì không trả nổi. nên huấn luyện \hl{$g_\theta$} — đoán $v(S)$, \yes{0 lượt llm}.}

  \[
    g_\theta\Bigl(q,\; S,\; \{\hat s(q,k)\}_{k \in S},\;
                  \{\tilde v_k\}_{k \in S},\; v(\emptyset)\Bigr)
    \;\longrightarrow\; \hat v(S)
  \]
  \dd{$\hat s$ = điểm xếp hạng, luôn có sẵn ·
      $\tilde v_k = v(\{k\}) - v(\emptyset)$ = đoạn $k$ một mình đóng góp bao nhiêu ·
      đảo thứ tự đoạn trong $S$ không đổi kết quả}

  \vspace{-0.1em}
  \begin{columns}[t]
  \begin{column}{0.49\textwidth}
    \begin{badb}[title={$\tau_q$ KHÔNG được làm đầu vào}]
      $g_\theta$ đoán \emph{điểm}; so với ngưỡng xảy ra sau, bên ngoài mô hình.
      \vspace{0.3em}

      lý do thật: \texttt{c2} chạy \emph{trước} \texttt{c4}, mà \texttt{c4} mới
      chốt $\tau_q$ — cho $\tau_q$ vào đầu vào là \no{tạo vòng tròn}.
      $\tau_q$ chỉ xuất hiện trong hàm mất mát lúc huấn luyện:
      \vspace{-0.2em}
      \[ \mathcal{L} = \mathcal{L}_{\text{cls}} + \lambda \mathcal{L}_{\text{rank}} \]
    \end{badb}
  \end{column}
  \begin{column}{0.47\textwidth}
    \begin{bx}[title={dữ liệu huấn luyện là nhật ký của chính các lượt đã gọi}]
      \begin{dia}
pha 1, bốn vòng           1.110.000 nhãn
verification pha 2        1.880.000
tỉa khi đáp số to           104.000
ca phải chạy lại đầy đủ     821.000
ca bị bỏ                     70.000
                        -----------
                          3.983.000
      \end{dia}
      \vspace{0.25em}
      \dd{không tốn thêm đồng nào: mỗi lần gọi llm để đo $v$, ta ghi lại cặp
          (tập con, điểm) — rồi dùng chính nó để dạy $g_\theta$}
    \end{bx}
  \end{column}
  \end{columns}
\end{frame}

\begin{frame}[fragile]{c2 · nghiệm thu $g_\theta$ bằng phép đo đúng}
  \framesubtitle{c2}
  \lead{dễ tự lừa ở đây. một con số đẹp vẫn có thể đi kèm một mô hình vô dụng.}

  \begin{columns}[t]
  \begin{column}{0.49\textwidth}
    \begin{badb}[title={cái bẫy: đo gộp mọi truy vấn}]
      tập to thì $v$ cao, tập nhỏ thì $v$ thấp — đúng với \emph{mọi} truy vấn.
      nên một mô hình \no{chỉ nhìn số đoạn} đã cho con số gộp rất đẹp.
      \vspace{0.3em}

      mà việc ta cần là so các tập \emph{trong cùng một} truy vấn, thường
      \emph{cùng cỡ}. ở đó mô hình đó phân biệt được \no{bằng 0} — mọi tập cỡ 3
      hòa điểm, và bước tìm thành chọn bừa.
    \end{badb}
  \end{column}
  \begin{column}{0.47\textwidth}
    \begin{keyb}[title={nên đo thế này}]
      \hl{auc tính trong từng truy vấn}, rồi mới lấy trung bình qua các truy vấn.
      \vspace{0.3em}

      \hl{sàn phải vượt:} đúng cái mô hình chỉ nhìn $|S|$ ở trên. không vượt nó
      thì $g_\theta$ chưa học được gì.
    \end{keyb}
    \begin{bx}[title={và thứ tự thử kiến trúc}]
      cây tăng cường (gbm) trên đặc trưng làm tay chạy \hl{trước}. mô hình
      neural phức tạp chỉ đáng chọn nếu vượt được nó trên phép đo trên.
    \end{bx}
  \end{column}
  \end{columns}
\end{frame}

\begin{frame}[fragile]{c5 · tìm đáp số — đi lên từ các tập phủ đủ mảnh}
  \framesubtitle{c5}
  \lead{không duyệt $2^{40}$ tập. chỉ duyệt những tập \emph{đã phủ hết mảnh},
        xếp theo cỡ từ nhỏ lên — và danh sách đó \texttt{c0} đã đưa sẵn.}

  \begin{columns}[t]
  \begin{column}{0.49\textwidth}
    \begin{bx}[title={các tầng}]
      \begin{dia}
tầng 0   các tập phủ đủ, nhỏ nhất
         = 6 ứng viên c0 đã liệt kê
         {3,12} {3,23} {3,31} {7,12} ...
tầng 1   một tập phủ đủ  +  1 đoạn
tầng 2   một tập phủ đủ  +  2 đoạn
...

trong mỗi tầng:
   chấm toàn bộ bằng g_theta   (0 lượt)
   giữ lại vài ứng viên đầu bảng
   đo v THẬT theo đúng thứ tự đó
      \end{dia}
    \end{bx}
  \end{column}
  \begin{column}{0.47\textwidth}
    \begin{keyb}[title={vì sao bỏ qua phần còn lại là an toàn}]
      một tập \emph{không} phủ hết mảnh mà vẫn vượt $\tau_q$ thì chỉ có thể vì
      \no{mồi tham số} — mà ta đã quyết loại loại ca đó.
      \vspace{0.3em}

      tìm ở đó là tìm đúng thứ mình sẽ vứt đi.
    \end{keyb}
    \begin{bx}[title={hai thứ tự nhiên đi kèm}]
      \textbf{1}\; mọi ứng viên đều đã đủ bằng chứng — không còn ca đoạn 19.
      \vspace{0.25em}

      \textbf{2}\; duyệt theo cỡ tăng dần, nên tập nhỏ luôn được thử trước:
      \yes{«nhỏ nhất trước» là miễn phí}, không cần một bước tỉa riêng.
    \end{bx}
  \end{column}
  \end{columns}
\end{frame}
